\section{Agentic Coding 闭环：训练对象到底是什么}

讲者给出了一张非常关键的系统图（coding-agent loop）：用户 issue 与代码仓库输入到模型，模型通过 tool calls 与环境模拟器交互，执行测试并迭代 patch，最后产出 PR 并接受下一轮反馈。这说明训练对象不是单次回答，而是一个 \textbf{多回合行动轨迹（trajectory）}。

这张图的工程含义有三层。第一，模型必须维护跨步状态，不能只做局部 token-level 最优。第二，环境接口（例如单元测试、仓库检索、命令执行）是训练信号的一部分。第三，用户反馈会形成持续约束，系统需要在后续轮次中修复此前失败。

\begin{knowledgebox}{Agent 训练中 ``环境'' 的角色}
在普通 SFT 任务里，数据通常是静态样本；在 agent 训练里，环境是动态数据生成器。每次 rollout 都会产生新的状态与新错误，进而改变下一轮训练分布。
\end{knowledgebox}

讲者提到实际系统中 CPU 资源规模可能是 GPU 的数倍到十倍，因为环境模拟与验证步骤（运行脚本、执行测试、沙箱交互）大多在 CPU 路径。换言之，Agent RL 是典型的 ``GPU + CPU 协同瓶颈'' 问题。

\begin{warningbox}{常见误解：把 Agent 训练简化为 ``工具标签学习''}
仅让模型学会输出 tool schema 并不等于具备 agent 能力。真正困难点在于：何时调用、如何回退、如何在失败后改写计划，以及如何在成本约束下终止探索。
\end{warningbox}

\begin{importantbox}{闭环训练的可迁移抽象}
\begin{itemize}
    \item 任务输入：issue / 指令 / 约束；
    \item 状态扩展：仓库、日志、历史对话；
    \item 行动集合：读写文件、运行测试、调用外部工具；
    \item 反馈信号：pass/fail、格式得分、长度成本、用户偏好。
\end{itemize}
这个抽象不仅适用于 coding，也适用于 data analysis、research assistant、workflow automation。
\end{importantbox}

\subsection{本章小结}
Agentic 训练对象是 ``环境中可执行的多步轨迹''，而不是单句回答。只有把环境反馈纳入训练回路，模型才会学会稳定完成任务而非偶然命中答案。

